\begin{statement}[C.8.6; Proved: the Mahler-measure identities]
Let $\Lambda\subset\C$ be a lattice of area $A$, let $v$ be a zero
of $\wp_\Lambda$, and let $d\mu=A^{-1}dA_z$ be Haar probability
measure on $\C/\Lambda$. Put $U(z)=\log N(z)$ and
$L=U(v)$. Then
\begin{align}
 6L&=\int_{\C/\Lambda}\log\left|\frac{\wp'(z)^2}{\wp(z)^3}\right|d\mu(z),
 \label{eq:ca-mahler}\\
 N(v)^2&=|\Delta(\Lambda)|^{1/6}
      \exp\left(-\int_{\C/\Lambda}\log|\wp(z)|d\mu(z)\right).
 \label{eq:ca-wpmean}
\end{align}
If $C(u)=1-(g_2/4)u^2-(g_3/4)u^3$, the pushforward of $d\mu$
under $u=1/\wp(z)$ is
\begin{equation}\label{eq:ca-nu}
 d\nu(u)=\frac1{2A}\frac{dA_u}{|u|\,|C(u)|}.
\end{equation}
In particular its unnormalized integral is exactly $2A$, and
\[
 N(v)^6=4\exp\left(\int_{\C}\log|C(u)|d\nu(u)\right).
\]
Writing $e_1,e_2,e_3$ for the three branch values and
$p(x)=\int\log|x-\wp(z)|d\mu(z)$, the maximum theorem is equivalently
\[
 \int\log|C|d\nu\leq\log2,\qquad
 p(0)\geq\frac{p(e_1)+p(e_2)+p(e_3)}3-\frac{\log2}3.
\]
Equality in either inequality occurs exactly when $g_3=0$.
\end{statement}
\begin{proof}
The mean of $U$ is zero by the theta product calculation in C.1,
and near zero
$U(z)=\log|z|+\frac1{12}\log|\Delta|+o(1)$.
The divisor of $\wp$ is $(v)+(-v)-2(O)$, with multiplicity
two when $v=-v$. Subtract its logarithmic divisor from the
corresponding combination of $U$'s. The result is harmonic
on the compact torus and hence constant. Its expansion at
zero gives
\[
 \log|\wp(z)|=U(z-v)+U(z+v)-2U(z)
                  +\frac16\log|\Delta|-2L.
\]
Integration proves \eqref{eq:ca-wpmean}.
Similarly the divisor of $\wp'$ is the sum of the three
nonzero half-periods minus $3(O)$. The multiplication-by-two
identity of C.1 gives the product of $N$ at those half-periods
as $2$. Since $\wp'(z)\sim-2z^{-3}$, comparison at zero yields
\[
 \log|\wp'(z)|=\sum_{w\in E[2]\setminus\{O\}}U(z-w)-3U(z)
                         +\frac14\log|\Delta|.
\]
Its mean is $\frac14\log|\Delta|$. Twice this identity minus
three times the preceding mean proves \eqref{eq:ca-mahler}.
All logarithmic singularities are locally integrable, so these
integrations also cover a zero at a half-period.

The map $z\mapsto x=\wp(z)$ has degree two. Its two inverse
sheets each have Jacobian $|\wp'(z)|^{-2}$; branch values form
a measure-zero set. Thus its pushforward density is
$2/(A|4x^3-g_2x-g_3|)$. Under $x=1/u$ the Jacobian is
$|u|^{-4}$ and the denominator is $4|u|^{-3}|C(u)|$.
This proves \eqref{eq:ca-nu} and its normalization. The identity
$\wp'^2/\wp^3=4C(u)$ gives the next formula.

Replacing $v$ above by a half-period gives
$p(e_j)=\frac16\log|\Delta|-2U(w_j)$; summing and using
$\sum_jU(w_j)=\log2$ proves the centroid formulation.
The equality characterization is that of C.8.5,
since $j=1728$ is equivalent to $g_3=0$ for a nonsingular lattice.
\end{proof}
